\section{Data and warning protocol}
\label{sec:methods}

\subsection{Simulated engine data and target}

The experiment uses the labeled training partition of the public
\texttt{tiny-N-CMAPSS} repository at commit
\texttt{b915997ff8d5}~\cite{tiny_ncmapss}; the reproducibility manifest records
the complete commit hash.
The source URL is \url{https://github.com/alovberg/tiny-N-CMAPSS}.
This repository documents a $100\times$-downsampled N-CMAPSS challenge subset.
The checked source files contain 90 simulated run-to-failure trajectories and
633,875 retained readings. We aggregate them to 6,825 engine-cycle rows, or
92.9 retained readings per cycle on average. The separate test partition has
no RUL labels and is not used here.

The target marks an engine cycle as inside a stated final maintenance window:
\begin{equation}
y=\begin{cases}
1, & \mathrm{RUL}\leq20,\\
0, & \mathrm{RUL}>20.
\end{cases}
\label{eq:maintenance-window}
\end{equation}
There are 1,890 positive cycle rows, a 27.7\% base rate. The value 20 is an
experimental warning boundary, not a recommended service interval. For each
trajectory, cycle plus RUL is constant; total simulated lifetime ranges from
48 to 100 cycles (mean 75.8, standard deviation 14.1). This variation makes
elapsed cycle a useful but incomplete baseline.

\subsection{Feature provenance and input sets}

N-CMAPSS distinguishes measured physical properties from auxiliary fields.
The primary measurement source used here is the 14-channel \texttt{Xs} group.
The two fields \texttt{Fc} (flight class) and \texttt{hs} (health state) are
auxiliary data and are excluded from every model. In particular, \texttt{hs}
is not treated as a sensor measurement. The experiment also excludes unit ID
and RUL from every model. Table~\ref{tab:feature-provenance} lists the exact
variables.

\begin{table}[H]
\centering
\caption{Input provenance. Means and population standard deviations are
computed within each engine cycle for every channel except elapsed cycle.}
\label{tab:feature-provenance}
\small
\begin{tabular}{p{0.24\textwidth}c p{0.60\textwidth}}
\toprule
Group & Count & Fields and use \\
\midrule
Measured physical (\texttt{Xs}) & 14 & \texttt{T24, T30, T48, T50, P15, P2, P21, P24, Ps30, P40, P50, Nf, Nc, Wf}; included in telemetry variants \\
Operating conditions (\texttt{W}) & 4 & \texttt{alt, Mach, TRA, T2}; used in cycle-plus-operating and all-observable inputs \\
Elapsed cycle & 1 & \texttt{cycle}; included only in variants that name it \\
Auxiliary fields & 2 & \texttt{Fc, hs}; excluded from every model \\
Identifier and label & 2 & \texttt{unit, RUL}; excluded from every model \\
\bottomrule
\end{tabular}
\end{table}

We compare four inputs with one fixed random-forest reference. The cycle-only
baseline has one feature. Cycle plus operating conditions has the cycle number
and eight within-cycle operating summaries. Measured telemetry only has 28
within-cycle summaries from the 14 measured channels. All observable inputs has
37 features: cycle, operating summaries, and measured-telemetry summaries. This
comparison separates three questions: how much a clock predicts, whether
operating state adds information, and whether measured telemetry reduces misses
beyond both.

Each variant uses a \texttt{RandomForestClassifier} with 250 trees, maximum
depth 12, minimum leaf size 3, balanced classes, and a primary random seed of
2026. The model is a fixed screening model. The study does not claim that it is
the best architecture for this data.

\subsection{Nested engine-held-out validation}

We use five outer \texttt{GroupKFold} splits grouped by engine. Each outer
test fold has 18 complete trajectories and the remaining 72 engines form the
outer training data. Every engine appears in exactly one outer test fold. Its
rows are absent from fitting and threshold selection for that fold.

Within each outer training set, four inner engine-grouped folds generate
out-of-fold probabilities. We choose the lowest probability threshold whose
pre-window false-alert rate is at or below 5\%. A final model then fits the 72
outer-training engines and predicts the 18 held-out engines. The pooled result
contains only these outer-test predictions. This procedure evaluates a warning
policy without choosing its threshold from test-engine labels.

\subsection{Nested model-family selection}

The fixed random forest is used for the information comparison so that the
input-set result does not depend on a particular architecture. We also evaluate
one secondary, all-observable warning policy that selects a model family within
each outer-training split. Its four candidates are the fixed random-forest
reference, an extra-trees classifier, a histogram-gradient-boosting classifier,
and a standardized RBF support-vector classifier. Their configurations are
recorded in the evaluation output.

For each outer fold, every candidate produces inner engine-grouped out-of-fold
probabilities. For each candidate, the inner probabilities determine the same
5\%-budget threshold as above. The selected candidate is the one with the
highest inner final-window recall at that threshold; ties break in order of
AUROC, average precision, precision, and lower false-alert rate. We then fit
only that selected candidate on all 72 outer-training engines and score the 18
held-out engines. Thus neither model-family selection nor threshold selection
uses labels from the outer test engines.

\subsection{Metrics, robustness, and warning timing}

We report AUROC, precision, final-window recall, and the pre-window false-alert
rate. We calculate 95\% percentile intervals by resampling complete engine
trajectories 2,000 times. We also report a fixed 0.5-threshold check on the
same outer-test probabilities. It is a diagnostic, not the policy used for the
primary result.

To test stochastic variation, we repeat the full fixed-random-forest input
comparison and the full selected-policy procedure with seeds 2024 through 2028.
The outer engine splits remain fixed, so these ranges test run-to-run variation
on the same data rather than a new-dataset replication.

For a first-warning analysis, let $c_u^{\mathrm{window}}$ be the first cycle
of engine $u$ with $\mathrm{RUL}\leq20$, and let $c_u^{\mathrm{warn}}$ be the
first cycle whose score meets its outer-fold threshold. The lead is
\begin{equation}
\ell_u=c_u^{\mathrm{window}}-c_u^{\mathrm{warn}}.
\label{eq:lead}
\end{equation}
Positive values are warnings before the window; negative values are warnings
after it begins. This timing is a property of the tested screening policy, not
a maintenance command.

\subsection{Reproducibility materials}

The project repository is
\url{https://github.com/prahaladtalur/N-CMAPSS-Engine-Prediction}. The
evaluation records the exact tiny-N-CMAPSS commit, SHA256 hashes of the two
downloaded data files, outer-test predictions, selected thresholds, and the
scripts used to produce the tables and figures. These materials identify the
corrected input definition and make the reported values reproducible.
